\lysh{22065_SOLUTION}{1}
\selectlanguage{greek}
\textbf{ΛΥΣΗ}

α) Το σημείο $A(0,0)$ είναι η αρχή των αξόνων. Επομένως,
\[ \vec{A\Delta} = (2,4) \text{ και } \vec{B\Gamma} = \vec{A\Gamma} - \vec{AB} = (10,4) - (8,0) = (2,4), \text{ οπότε } \vec{A\Delta} = \vec{B\Gamma}, \]
συνεπώς το τετράπλευρο $AB\Gamma\Delta$ είναι παραλληλόγραμμο.

β) Αν $M$ και $N$ είναι τα μέσα των πλευρών $AB$ και $B\Gamma$ αντίστοιχα, τότε
\[ M\left(\frac{x_A+x_B}{2}, \frac{y_A+y_B}{2}\right) \text{ ή ισοδύναμα } M\left(\frac{0+8}{2}, \frac{0+0}{2}\right), \text{ δηλαδή } M(4,0). \]
\[ N\left(\frac{x_B+x_\Gamma}{2}, \frac{y_B+y_\Gamma}{2}\right) \text{ ή ισοδύναμα } N\left(\frac{8+10}{2}, \frac{0+4}{2}\right), \text{ δηλαδή } N(9,2). \]

γ) Εξίσωση της ευθείας $A\Gamma$:
\[ \frac{y-y_A}{x-x_A} = \frac{y_\Gamma-y_A}{x_\Gamma-x_A} \Leftrightarrow \frac{y-0}{x-0} = \frac{4-0}{10-0} \Leftrightarrow y = \frac{4}{10}x \Leftrightarrow y = \frac{2}{5}x \quad (1). \]
Εξίσωση της ευθείας $\Delta M$:
\[ \frac{y-y_M}{x-x_M} = \frac{y_\Delta-y_M}{x_\Delta-x_M} \Leftrightarrow \frac{y-0}{x-4} = \frac{4-0}{2-4} \Leftrightarrow y = \frac{4}{-2}(x-4) \Leftrightarrow y = -2(x-4) \Leftrightarrow y = -2x+8 \quad (2). \]
Εξίσωση της ευθείας $\Delta N$:
\[ \frac{y-y_N}{x-x_N} = \frac{y_\Delta-y_N}{x_\Delta-x_N} \Leftrightarrow \frac{y-2}{x-9} = \frac{4-2}{2-9} \Leftrightarrow y-2 = \frac{2}{-7}(x-9) \Leftrightarrow y = -\frac{2}{7}x + \frac{18}{7} + 2 \Leftrightarrow y = -\frac{2}{7}x + \frac{32}{7} \quad (3). \]
Το σημείο $K$ είναι σημείο τομής των ευθειών $A\Gamma$ και $\Delta M$, συνεπώς οι συντεταγμένες του είναι λύση του συστήματος των εξισώσεων (1) και (2): $y = \frac{2}{5}x$ και $y = -2x+8$.
Εξισώνοντας τα δεξιά μέλη έχουμε $ \frac{2}{5}x = -2x+8 $ δηλαδή $ x = \frac{10}{3} $, άρα $ K\left(\frac{10}{3}, \frac{4}{3}\right) $.

Το σημείο $\Lambda$ είναι σημείο τομής των ευθειών $A\Gamma$ και $\Delta N$, συνεπώς οι συντεταγμένες του είναι λύση του συστήματος των εξισώσεων (1) και (3): $y = \frac{2}{5}x$ και $y = -\frac{2}{7}x + \frac{32}{7}$.
Εξισώνοντας τα δεξιά μέλη έχουμε $ \frac{2}{5}x = -\frac{2}{7}x + \frac{32}{7} $ δηλαδή $ x = \frac{20}{3} $, άρα $ \Lambda\left(\frac{20}{3}, \frac{8}{3}\right) $.

δ) Το σημείο $A(0,0)$ είναι η αρχή των αξόνων και έχουμε:
Από την εκφώνηση: $ \vec{A\Gamma} = (10,4) $
Από την απάντηση του γ) ερωτήματος: $ \vec{AK} = \left(\frac{10}{3}, \frac{4}{3}\right) $ και $ \vec{A\Lambda} = \left(\frac{20}{3}, \frac{8}{3}\right) $.
Άρα, $ \vec{AK} = \frac{1}{3}\vec{A\Gamma} $ και $ \vec{A\Lambda} = \frac{2}{3}\vec{A\Gamma} $.

% TODO: TikZ σχήμα: Ευθύγραμμο τμήμα A-Γ με σημεία K, Λ που το τριχοτομούν. (A -- K -- Λ -- Γ)

Συνεπώς, τα σημεία $K$ και $\Lambda$ τριχοτομούν τη διαγώνιο $A\Gamma$.
\endlysh